\subsection{Closed evaluation of the spanning matrix}
\label{sec:comp-supp-box-tensor}
The finite two-parameter terms in
\eqref{eq:comp-supp-spanning-integrated} can be integrated without
constructing a separate primitive for each numerator coefficient.
The numerator is a polynomial of degree three in the loop momentum.
Differences of propagator denominators therefore reduce its tensor
integrals to a small set of scalar integrals. We derive this reduction
before listing the scalar answers. All momenta remain in units of the
electron mass.

Write the four denominators as
\begin{equation}
 d_i=(\ell-q_i)^2-m_i^2+\ii0,\qquad
 (q_0,q_1,q_2,q_3)=(0,p,Q,p'),\qquad
 (m_0^2,m_1^2,m_2^2,m_3^2)=(0,1,1,1).
\end{equation}
The photon mass is kept in the infrared-divergent scalar integrals.
Setting it to zero in the algebraic coefficients omits only terms
that vanish as $\lambda\to0$. For an ordered subset $I$ of denominator
indices define
\begin{equation}
 T_I^{\mu_1\cdots\mu_r}
 =\frac{e^{\gamma_E\epsilon}\mu^{2\epsilon}}{\ii\pi^{d/2}}
   \int\dd^d\ell\,
   \frac{\ell^{\mu_1}\cdots\ell^{\mu_r}}
        {\prod_{i\in I}d_i},\qquad d=4-2\epsilon.
 \label{eq:comp-supp-box-tensor-definition}
\end{equation}
The factor $e^{\gamma_E\epsilon}$ fixes the normalization so that
the common pole is $U=1/\epsilon$, without an accompanying
$-\gamma_E$. Since the full box is ultraviolet finite, this common
factor does not change its four-dimensional value. We choose
$\mu=m$ in the dimensionless finite parts below.

\paragraph{Contractions from denominator differences.}
Choose an index $a\in I$ as the origin and put
$L=\ell-q_a$. Let $v_i=q_i-q_a$ for $i\in I\setminus\{a\}$.
Define the Gram matrix, its dual vectors and the transverse projector
by
\begin{equation}
 G_{ij}=v_i\cdot v_j,\qquad
 w_i^\mu=\sum_j(G^{-1})_{ij}v_j^\mu,\qquad
 P^{\mu\nu}=\eta^{\mu\nu}-\sum_iw_i^\mu v_i^\nu.
 \label{eq:comp-supp-box-projector}
\end{equation}
These formulas first apply when the vectors $v_i$ are linearly
independent and $\det G\ne0$. Singular limits are taken in the
combined tensor, not in isolated inverse-Gram coefficients.
Thus $w_i\cdot v_j=\delta_{ij}$ and
$P^\mu{}_{\mu}=d-(|I|-1)$. The two required identities are
\begin{equation}
 2L\cdot v_i=d_a-d_i+f_i,\qquad
 L^2=d_a+m_a^2,\qquad
 f_i=v_i^2+m_a^2-m_i^2.
 \label{eq:comp-supp-box-denominator-differences}
\end{equation}
For the following formulas every tensor uses the numerator $L$.
A pinched tensor uses this same $L$, even when its denominator set
no longer contains $a$. This convention avoids an implicit shift
of the numerator in the reduction.

Suppressing the set $I$, the first identity gives
\begin{equation}
 V_i^{\mu_2\cdots\mu_r}
 \equiv v_{i\mu_1}T^{\mu_1\cdots\mu_r}
 =\frac12\left[
 T_{I\setminus\{a\}}^{\mu_2\cdots\mu_r}
 -T_{I\setminus\{i\}}^{\mu_2\cdots\mu_r}
 +f_iT_I^{\mu_2\cdots\mu_r}\right].
 \label{eq:comp-supp-box-pinch-recurrence}
\end{equation}
The loop rank falls by one on the right-hand side. Define
\begin{equation}
 H_i=V_i,\qquad
 H_{ij}=v_{j\nu}V_i^\nu,\qquad
 H_{ijk}=v_{j\nu}v_{k\rho}V_i^{\nu\rho},\qquad
 n=|I|-1.
\end{equation}
Rotations in the space transverse to all $v_i$ then determine
ranks one through three completely,
\begin{align}
 T^\mu&=\sum_iw_i^\mu H_i,\\
 T^{\mu\nu}&=\sum_{ij}w_i^\mu w_j^\nu H_{ij}
                      +P^{\mu\nu}Z,\\
 Z&=\frac{T_{I\setminus\{a\}}+m_a^2T_I
                   -\sum_{ij}(G^{-1})_{ij}H_{ij}}{d-n},\\
 T^{\mu\nu\rho}&=\sum_{ijk}w_i^\mu w_j^\nu w_k^\rho H_{ijk}
       +\sum_i\left(P^{\mu\nu}w_i^\rho
                   +P^{\mu\rho}w_i^\nu
                   +P^{\nu\rho}w_i^\mu\right)Z_i,\\
 Z_i&=\frac{v_{i\mu}
       (T_{I\setminus\{a\}}^\mu+m_a^2T_I^\mu)
                 -\sum_{jk}(G^{-1})_{jk}H_{ijk}}{d-n}.
 \label{eq:comp-supp-box-rank-three}
\end{align}
The equations for $Z$ and $Z_i$ follow by taking a Lorentz trace
and using the second identity in
\eqref{eq:comp-supp-box-denominator-differences}.
To change the origin in a pinched integral, expand each factor
$L=\widetilde L+(q_b-q_a)$ in its numerator. This is a finite
binomial expansion of tensors of lower or equal rank.
Consequently a rank-three box requires only rank-two triangles,
rank-one bubbles and scalar tadpoles.

Although the original box is ultraviolet finite, some pinched
integrals are not. Their poles must be carried through the
recurrence before taking $d=4$. In particular,
\begin{equation}
 \frac{x_{-1}/\epsilon+x_0}{c-2\epsilon}
 =\frac{x_{-1}}{c\epsilon}+\frac{x_0}{c}
                          +\frac{2x_{-1}}{c^2}
                          +\mathcal O(\epsilon).
 \label{eq:comp-supp-box-rational-finite}
\end{equation}
The last term is needed for the finite box. This also specifies
how the calculation handles the rational terms without assuming
that a four-dimensional tensor reduction is sufficient.

\paragraph{Scalar integrals.}
Let $S=Q^2$, $D=S-1$, $T=-t\geq0$, and define
\begin{equation}
 a=\frac{\sqrt T}{2},\qquad h=\sqrt{1+a^2},\qquad
 r=\operatorname{arsinh}a,\qquad \kappa=\frac ah.
\end{equation}
For clarity, denote scalar integrals by
$A_i=T_{\{i\}}$, $B_{ij}=T_{\{i,j\}}$,
$C_{ijk}=T_{\{i,j,k\}}$ and $D_{0123}=T_{\{0,1,2,3\}}$.
The tadpoles and bubbles are
\begin{align}
 A_0&=0,\qquad A_1=A_2=A_3=U+1,\\
 B_{01}&=B_{03}=U+2,\qquad B_{12}=B_{23}=U,\\
 B_{02}&=U+2+\frac{1-S}{S}\ln(1-S-\ii0),\\
 B_{13}&=U+2-\frac{2h}{a}r.
 \label{eq:comp-supp-box-bubbles}
\end{align}
The continuous values at $S=0$ and $T=0$ are used.
These expressions follow by integrating
\begin{equation}
 B_{ij}=U-\int_0^1\dd x\,
 \ln\!\left[xm_i^2+(1-x)m_j^2
                    -x(1-x)(q_i-q_j)^2-\ii0\right].
\end{equation}
For equal masses the rank-one bubble is exactly
$(q_i+q_j)^\mu B_{ij}/2$. This formula is also valid when
$(q_i-q_j)^2=0$, where a one-vector Gram inverse would be
singular. For unequal masses use
\begin{equation}
 T_{\{i,j\}}^\mu=q_i^\mu B_{ij}
 +(q_j-q_i)^\mu
 \frac{A_j-A_i+[(q_j-q_i)^2+m_i^2-m_j^2]B_{ij}}
      {2(q_j-q_i)^2}.
\end{equation}
For $i=0$, $m_j=1$, the coefficient in the second term approaches
$U/2+1/4$ as $(q_j-q_0)^2\to0$.

The four triangles reduce to three distinct functions,
\begin{align}
 C_{012}=C_{023}
 &=-\frac{\operatorname{Li}_2(S+\ii0)-\pi^2/6}{S-1},\\
 C_{123}&=-\frac{2r^2}{T},\\
 C_{013}&=\frac r{ah}\ln\lambda-\frac12\mathcal H(T),\\
 \mathcal H(T)&=
 \frac{2r\ln(2h)-r^2-\operatorname{Li}_2(-e^{-2r})
                        +\operatorname{Li}_2(-1)}{ah}.
 \label{eq:comp-supp-box-triangles}
\end{align}
The overall minus signs in the finite triangles follow from
$\int\dd^d\ell/(\ii\pi^{d/2})$ with three Feynman denominators.
The first line is the previously integrated vertex moment $-J_{00}$.
For the pure-electron triangle, integrate its simplex denominator
$1+T xz$ to obtain $2r^2/T$ before applying this sign.
The infrared triangle follows by writing the sum of its two
electron parameters as $\rho$,
\begin{align}
 C_{013}&=-\int_0^1\dd z\int_0^1
 \frac{\rho\dd\rho}{\rho^2[1+Tz(1-z)]+(1-\rho)\lambda^2-\ii0}\\
 &=g_0(t)\ln\lambda-\frac12\int_0^1
          \frac{\ln[1+Tz(1-z)]}{1+Tz(1-z)}\dd z+o(1).
\end{align}
To evaluate the last integral, set $2z-1=y$ and then
$\kappa y=\tanh u$. Its upper boundary is $u=r$, giving
\begin{align}
 \mathcal H(T)&=\frac1{ah}
 \left[2r\ln h-2\int_0^r\ln\cosh u\,\dd u\right],\\
 \int_0^r\ln\cosh u\,\dd u
 &=\frac{r^2}{2}-r\ln2
  +\frac{\operatorname{Li}_2(-e^{-2r})-\operatorname{Li}_2(-1)}2.
\end{align}
This proves the last line of
\eqref{eq:comp-supp-box-triangles}. Its boundary values are
$C_{123}\to-1/2$ and $\mathcal H(T)\to0$ at $T=0$.
Finally the scalar box is the result already derived in
\eqref{eq:comp-supp-box-dilog},
\begin{equation}
 D_{0123}=\frac1D\left[
 \frac r{ah}\ln\frac{\lambda}{-D-\ii0}
 +\frac{\operatorname{Li}_2(\kappa)-\operatorname{Li}_2(-\kappa)}{2ah}
 \right]+o(1).
\end{equation}
No parameter integral remains in this list.

\paragraph{Physical boundary values.}
The parameter denominators have the boundary value
$\Delta-\ii0$ inherited from the original propagators.
For $S>1$ this gives
$\ln(1-S-\ii0)=\ln(S-1)-\ii\pi$.
Since $\dd\operatorname{Li}_2(S)/\dd S=-\ln(1-S)/S$,
the corresponding dilogarithm is approached from the upper
half-plane. This fixes its imaginary part by continuation from
$S<1$, without choosing a sign from a numerical comparison.
The functions of $T\geq0$ in the other lines are real.
In particular $-e^{-2r}\in[-1,0)$ and $\kappa\in[0,1)$,
so they do not cross a dilogarithm cut.

\paragraph{Assembly of the matrix.}
Expand the original numerator as
\begin{equation}
 N(\ell)=\sum_{r=0}^3
   \mathcal N_{\mu_1\cdots\mu_r}
           \ell^{\mu_1}\cdots\ell^{\mu_r}
 =\gamma^\alpha(\slashed p'-\slashed\ell+1)
 B(\slashed Q-\slashed\ell+1)A
 (\slashed p-\slashed\ell+1)\gamma_\alpha,
\end{equation}
where $A$ and $B$ are the incident and outgoing photon vertex
matrices in the chosen electron ordering. The complete closed
spanning matrix is then
\begin{equation}
 \boxed{\widehat W=\sum_{r=0}^3
       \mathcal N_{\mu_1\cdots\mu_r}
       \left.T_{\{0,1,2,3\}}^{\mu_1\cdots\mu_r}
       \right|_{\epsilon^0}.}
 \label{eq:comp-supp-box-closed-matrix}
\end{equation}
Equations~\eqref{eq:comp-supp-box-pinch-recurrence}--
\eqref{eq:comp-supp-box-triangles} give every coefficient through
finite algebra and logarithms or dilogarithms.
The ultraviolet poles cancel in each box tensor.
At exactly forward or backward kinematics the three-vector Gram
matrix is singular. The physical answer is the continuous limit
of the closed result at interior scattering angles.
This completes the loop-parameter integrations of the spanning
matrix. The subsequent spin interference and elastic-angle
integration are separate steps toward the total cross section.

\section{Virtual interference and angular integration}
\subsection{Closed interference coefficients for the box}
\label{sec:comp-supp-box-interference}
For the angular integration, introduce positive variables
$x=s-1$ and $y=1-u$, so $t=y-x$. Contract each external photon
pair with $-\eta_{\mu\nu}$ and average over the two initial electron
spins and two initial photon polarizations. Individual diagram
groups do not satisfy the photon Ward identity. The same covariant
polarization contraction must therefore be used for the box,
vertex and self-energy groups before they are added.

Let $\mathcal T_{\mathrm{box},s}$ denote the bilinear between
the $s$-ordered spanning matrix and the full Born matrix, with
$\alpha_{\mathrm T}/(4\pi)$ removed from the former. Define finite
scalar parts $\overline A_1=A_1-U$ and
$\overline B_{ij}=B_{ij}-U$. Direct Clifford traces and the preceding
reduction give
\begin{equation}
 \mathcal T_{\mathrm{box},s}
 =f_A\overline A_1+f_a\overline B_{01}
 +f_b\overline B_{02}+f_t\overline B_{13}
 +f_c C_{012}+f_d C_{013}+f_e C_{123}
 +f_DD_{0123}+f_0.
 \label{eq:comp-supp-box-interference-closed}
\end{equation}
The coefficients are real rational functions,
\begin{align}
 f_A&=\frac{4(xy-x+y)}{xy(x+1)},\\
 f_a&=-\frac{8P_a}{x^2y(x-y+4)},\\
 f_b&=\frac{4P_b}{x^2y(x+1)},\\
 f_t&=-\frac{4(x-y)(xy-x+y)}{xy(x-y+4)},\\
 f_c&=\frac{4(2x^2-2xy-2x+y^2-8)}{y},\\
 f_d&=-\frac{2(2x^3-4x^2y+2x^2+xy^2-6xy-12x+y^3-4y)}{xy},\\
 f_e&=-\frac{2(x-y)(2x^2-2xy-2x+y^2-2y-4)}{xy},\\
 f_D&=\frac{2P_D}{xy},\qquad
 f_0=\frac{4(xy-x+y)}{xy}.
\end{align}
Here the polynomials in the longer coefficients are
\begin{align}
 P_a={}&x^3y-2x^2y^2+6x^2y+xy^3-5xy^2-2xy
                +8x+y^3-2y^2-8y,\\
 P_b={}&3x^3y-x^3-2x^2y^2+3x^2y+4x^2
                -4xy^2-4xy+4x-2y^2-4y,\\
 P_D={}&2x^4-4x^3y+3x^2y^2+2x^2y-24x^2
                -xy^3-2xy^2+24xy-16x-8y^2+16y.
\end{align}
The $u$-ordered coefficient of each corresponding scalar integral
is obtained by the simultaneous replacement
$(x,y)\mapsto(-y,-x)$. In that ordering the off-shell invariant
in $\overline B_{02}$, $C_{012}$ and $D_{0123}$ is $S=1-y$.
This relation follows from exchanging the two photon insertions
in the original numerator. It also follows by independent trace
evaluation of that ordering.

The calculation of these coefficients uses only the Clifford
identity $\{\gamma^\mu,\gamma^\nu\}=2\eta^{\mu\nu}$ and the tensor
recurrences above. For example, the scalar trace numerator is
at most cubic in $\ell$. Each factor of $\ell^2$ cancels $d_0$,
and each contraction of $\ell$ with an external momentum is
reduced using \eqref{eq:comp-supp-box-denominator-differences}.
There are seventeen scalar monomials before this reduction in
each photon ordering. The coefficients multiplying $U$ cancel
exactly. Thus \eqref{eq:comp-supp-box-interference-closed} is an
explicit finite expression, including its infrared logarithm.

\subsection{A rational variable for the elastic angle}
For fixed $x>0$ the physical interval is
$x/(1+x)\leq y\leq x$. The substitution
\begin{equation}
 v=\sqrt{\frac{x-y}{x-y+4}},\qquad
 y=\frac{x-(x+4)v^2}{1-v^2},\qquad
 0\leq v\leq\frac{x}{x+2},\qquad
 -\dd y=\frac{8v}{(1-v^2)^2}\dd v
 \label{eq:comp-supp-box-angle-rationalization}
\end{equation}
removes the square root from every spacelike scalar integral.
In particular $r=\operatorname{artanh}v$ and
$ah=v/(1-v^2)$. This substitution follows directly from the
quadratic $T(T+4)$ in the scalar box, and its upper boundary follows
from the backward-scattering value $T=x^2/(1+x)$.

The generalized polylogarithms (GPLs) used here are defined recursively by
\begin{align}
 G(\varnothing;v)&=1,\\
 G(a_1,\ldots,a_n;v)
 &=\int_0^v\frac{\dd t}{t-a_1}G(a_2,\ldots,a_n;t),\\
 G(\underbrace{0,\ldots,0}_{n};v)&=\frac{\ln^n v}{n!}.
\end{align}
The last line specifies the logarithmically singular lower boundary
rather than an ordinary convergent integral. The ordered list
$(a_1,\ldots,a_n)$ is called a word, its entries are letters, and its
weight is its length $n$. The alphabet is the finite set of letters
needed by a given expression. These terms describe the iterated
integrals, not additional physical assumptions.
For a real letter outside the integration interval,
$G(a;v)=\ln(1-v/a)$. Letters on the path inherit their displacement
from the original Feynman denominators, as specified for the crossed
triangle below.

At a logarithmic endpoint, a regularized value means the constant
term after explicitly separating powers of the endpoint logarithm.
This choice is also called a tangential base point. At the origin
it is fixed by the all-zero formula above and by the vanishing limit
of words containing a nonzero letter. A full convergent combination
is formed before taking an endpoint limit. Define
\begin{equation}
 R(v)=\frac{G(-1;v)-G(1;v)}2,
 \qquad Z=\sqrt{\frac{x}{x+4}}.
\end{equation}
Then the needed logarithms and spacelike functions have the
explicit word representations
\begin{align}
 \ln y&=\ln x+G(Z;v)+G(-Z;v)-G(1;v)-G(-1;v),\\
 g_0(t)&=\frac{1-v^2}{v}R(v),\\
 g_a(t)&=\frac{1-v^2}{2v}\left[G(0,-1;v)-G(0,1;v)\right],\\
 \mathcal H(T)&=\frac{1-v^2}{2v}
 \left[G(1,1;v)-G(1,-1;v)+G(-1,1;v)-G(-1,-1;v)\right],\\
 C_{123}&=-\frac{1-v^2}{2v^2}R(v)^2,\qquad
 \overline B_{13}=2-\frac{2R(v)}v.
\end{align}
Products of words are reduced by interleaving their letters while
preserving the order within each word. This identity follows by
partitioning the product of their ordered integration domains.
For example $G(a)G(b)=G(a,b)+G(b,a)$.
These relations turn the closed box interference into rational
functions of $v$ multiplying words of weight at most two. They
provide the starting point for its remaining elastic-angle integral.

\subsection{Integrating the crossed rational subgroup first}
\label{sec:comp-crossed-first}
Before integrating the crossed box in the rationalized angle, separate
the scalar terms whose coefficients and function arguments are rational
in the original invariant $y$.
Transforming those terms to the angle variable would introduce an auxiliary root $\rho_u=\sqrt{(x-1)/(x+3)}$, whose
apparent threshold at $x=1$ is absent from the original scalar
integral. This root is introduced by applying the spacelike
variable $v$ to a subgroup that is already rational in the
original invariant $y$.
This subgroup can be integrated before that change of variable.
Here $x=s-1>0$ and $x/(1+x)\leq y\leq x$, in the mass units
used for the box calculation.

For this subgroup write $A=\overline A_1$, $B_s=\overline B_{02}$,
and $C_s=C_{012}$, evaluated at the crossed invariant $S=1-y$.
The subscript $s$ in these local scalar names does not label the
photon ordering. Include also the explicit constant in the crossed-box
trace. The scalar functions are
\begin{equation}
 A=1,\qquad B_s=2+\frac{y\ln y}{1-y},\qquad
 C_s=\frac{\operatorname{Li}_2(1-y)-\pi^2/6}{y}.
\end{equation}
Define the two GPLs
\begin{equation}
 F(y)=G(1,0;y)=\frac{\pi^2}{6}-\operatorname{Li}_2(1-y),
 \qquad W(y)=G(0,1,0;y).
\end{equation}
Their differential equations are
$F'(y)=\ln y/(y-1)$ and $W'(y)=F(y)/y$.
In particular, $F$ is real and analytic for every $y>0$.
The apparent pole of its derivative at $y=1$ is removable.
These statements fix the continuation of the GPL combination
without assigning separate imaginary parts at an artificial pole.

The selected crossed-box integrand is
\begin{equation}
 I_{\mathrm{rat}}(x,y)=R(x,y)+Q(x,y)\ln y-D(x,y)F(y),
\end{equation}
where exact partial fractions of the trace coefficients give
\begin{align}
 R={}&\frac{16(x-2)}{y^2}+\frac{28x+12}{x}
       +\frac4{x(y-1)}-\frac{8(2x^2+x-4)}{xy},\\
 Q={}&\frac{8(x-2)}y+\frac{4(x-6)}{x(y-1)}
       -\frac{12x+4}{x}-\frac4{x(y-1)^2},\\
 D={}&-\frac{8y}x+\frac{8x-8}x-\frac{4(x^2-8)}{xy}.
\end{align}
The only apparent double pole is integrated with its accompanying
simple pole. The exact identity
\begin{equation}
 \frac1{y-1}-\frac{\ln y}{(y-1)^2}
 =\frac{\dd}{\dd y}\frac{\ln y}{y-1}+\frac1y
\end{equation}
shows explicitly why no isolated $\ln(1-y)$ is needed.
Integration by parts of the two polynomial coefficients of $F$
then gives the finite primitive
\begin{align}
 P(x,y)={}&-\frac{16(x-2)}y+\frac{(32x+28)y+y^2}{x}\\
 &+\frac{-16x^2-8x+36-(4x+16)y-2y^2}{x}\ln y
       +4(x-2)\ln^2y\\
 &+\frac4x\frac{\ln y}{y-1}
       +\frac4x[y^2+(2-2x)y+3x-9]F(y)\\
 &+\frac{4(x^2-8)}xW(y).
\end{align}
Direct differentiation, using the two displayed GPL differential
equations, gives $\partial_yP=I_{\mathrm{rat}}$ as a rational
identity in the independent symbols $\ln y,F,W$.
The contribution to the unnormalized crossed angular integral is
therefore exactly
\begin{equation}
 J_{u,\mathrm{rat}}(x)=P(x,x)-P\left(x,\frac{x}{1+x}\right).
\end{equation}
Since $\dd c=2s\,\dd y/x^2$ and the interference contributes
$(4s)^{-1}\int\operatorname{Re}\mathcal T\,\dd c$, this subgroup
enters the cross-section coefficient as $J_{u,\mathrm{rat}}/(2x^2)$.
At $x=1$ the upper endpoint quotient is assigned
its analytic value $\ln y/(y-1)=1$, so the same expression covers
the apparent root collision there.

The lower endpoint introduces no new alphabet. Write
$f(x)=x/(1+x)$. The first identity below follows directly from
the logarithm of this quotient. For the other two, differentiation
and the boundary values $F(0)=W(0)=0$ prove
\begin{align}
 \ln f(x)&=G(0;x)-G(-1;x),\\
 F(f(x))&=-G(-1,0;x)+G(-1,-1;x),\\
 W(f(x))&=-G(0,-1,0;x)+G(0,-1,-1;x)\\
        &\quad+G(-1,-1,0;x)-G(-1,-1,-1;x).
\end{align}
Thus this entire subgroup has constant letters $\{-1,0,1\}$.
The remaining crossed masters are
$B_a=\overline B_{01}$, $B_t=\overline B_{13}$,
$C_t=C_{013}$, $C_m=C_{123}$, and $D_b=D_{0123}$.
Their rational trace coefficients have denominators built only
from $x$, $y$, and $x-y+4$. In particular none has a pole at
$y=1$. After the spacelike substitution their words involve
$0,\pm1,\pm\sqrt{x/(x+4)}$ and do not involve $\rho_u$.
Consequently $\rho_u$ disappears within the crossed box itself,
before adding soft or real-radiation contributions. This removes an auxiliary root within the virtual calculation.

There is also no remaining positive interior letter in the angular
remainder. If $V=x/(x+2)$ and $z=\sqrt{x/(x+4)}$, then
\begin{equation}
 z^2-V^2=\frac{4x}{(x+4)(x+2)^2}>0.
\end{equation}
Thus $0<V<z<1$. The remainder can be integrated along the real
$v$ interval without detouring around an artificial crossed root.

To fix the lower boundary in the subsequent function-basis
conversion, the series at positive $x\to0$ follows from
\begin{align}
 F(y)&=\sum_{n=1}^{\infty}y^n
             \left(\frac1{n^2}-\frac{\ln y}{n}\right),\\
 W(y)&=\sum_{n=1}^{\infty}y^n
             \left(\frac2{n^3}-\frac{\ln y}{n^2}\right).
\end{align}
They give
\begin{align}
 J_{u,\mathrm{rat}}(x)
 ={}&x(32\ln x-32)+x^2\left(\frac{92}{3}-24\ln x\right)\\
 &+x^3\left(\frac{88}{3}\ln x-\frac{236}{9}\right)
             +O(x^4\ln x).
\end{align}
In particular this unnormalized subgroup vanishes at the lower
boundary. The series fixes the boundary constant from the derived
integrand itself, without numerical fitting or a literature value.


\subsection{Angular primitives and endpoint assembly}
The preferred construction applies the finite primitive rules below to
the direct box and to the seventeen-word crossed remainder, adding the
rational subgroup boundary obtained above. For an independent algebraic
representation, one may instead transform the entire crossed integrand.
This introduces twenty-nine words. We give its continuation explicitly
so that both representations have a fully specified physical branch.

The crossed real-photon triangle is also a weight-two expression.
Put $\eta^2=(x-1)/(x+3)$ and define coefficients
\begin{equation}
 (c_Z,c_{-Z},c_1,c_{-1})=(1,1,-1,-1),\qquad
 (d_\eta,d_{-\eta},d_1,d_{-1})=(-1,-1,1,1).
\end{equation}
Differentiating its dilogarithm gives
\begin{equation}
 \frac{\dd}{\dd v}\operatorname{Li}_2(1-y(v)+\ii0)
 =\frac{y'(v)}{1-y(v)+\ii0}\ln(y(v)-\ii0),\qquad
 \frac{y'}{1-y}=\sum_{a\in\{\eta,-\eta,1,-1\}}\frac{d_a}{v-a}.
\end{equation}
Integrating each simple kernel therefore gives the explicit result
\begin{equation}
 \begin{aligned}
 \operatorname{Li}_2(1-y(v)+\ii0)-\frac{\pi^2}{6}
 ={}&\operatorname{Li}_2(1-x)-\frac{\pi^2}{6}\\
 &+\sum_a d_a\left[\ln x\,G(a;v)
                +\sum_b c_bG(a,b;v)\right].
 \end{aligned}
 \label{eq:comp-supp-box-crossed-triangle-words}
\end{equation}
The sums have exactly the four entries listed above.
For $x>1$ the root $\eta>0$ lies inside the angular interval.
Indeed, with $v_{\max}=x/(x+2)$,
\begin{equation}
 v_{\max}^2-\eta^2=\frac{4}{(x+2)^2(x+3)}>0,\qquad
 Z^2-v_{\max}^2=\frac{4x}{(x+4)(x+2)^2}>0.
\end{equation}
Thus the only positive auxiliary root crossed before the endpoint
is $\eta$, while the zero $Z$ of $y$ is outside the path.
It is an intermediate zero of $1-y$, not a singularity of the
triangle: $\ln y$ vanishes there. Its displacement follows from
$1-y+\ii0=0$. Since $y'(\eta)<0$, the positive root is below
the real axis and the negative root is above it. More explicitly,
\begin{equation}
 y'(v)=-\frac{8v}{(1-v^2)^2},\qquad
 v_{\rm pole}=\eta+\frac{\ii\varepsilon}{y'(\eta)}
                   +O(\varepsilon^2),\qquad \varepsilon>0.
\end{equation}
Here $\varepsilon$ is the infinitesimal Feynman displacement, not
a soft-energy separator. The GPL terms
in \eqref{eq:comp-supp-box-crossed-triangle-words} use those
boundary values together. Their apparent singularity cancels.
For $x<1$ the two roots are off the real integration interval.
The value at $x=1$ is obtained by continuity of the combined
expression. The reality of the combined result can also be proved
without evaluating any GPL. Throughout the path $y>0$, hence
$1-y<1$ is on the real branch of the dilogarithm. At $y=1$,
$\ln y/(1-y)\to-1$, so its displayed derivative is regular.
That derivative and the real initial value at $v=0$ uniquely fix
a real function on the full interval. Consequently the imaginary
parts introduced by the separated GPL terms must cancel.

\paragraph{Finite integration of the words.}
The remaining rational coefficients need no new special functions.
For a rational function $F(v)$, separate its polynomial and pole
parts,
\begin{equation}
 F(v)=\sum_{k\geq0}p_kv^k
       +\sum_a\sum_{j=1}^{m_a}\frac{r_{aj}}{(v-a)^j}
       =P'(v)+\sum_a\frac{r_{a1}}{v-a},
\end{equation}
where all sums are finite and
\begin{equation}
 P(v)=\sum_{k\geq0}\frac{p_kv^{k+1}}{k+1}
       -\sum_a\sum_{j=2}^{m_a}
                    \frac{r_{aj}}{(j-1)(v-a)^{j-1}}.
\end{equation}
The coefficients $r_{aj}$ are fixed algebraically by derivatives
of $(v-a)^{m_a}F(v)$ at $v=a$. This form is the rational
Hermite reduction, written here as partial fractions.
It gives the following finite primitive for the empty word and
for words of lengths one and two,
\begin{align}
 \Phi_{\varnothing}[F](v)
   &=P(v)+\sum_\alpha r_{\alpha1}G(\alpha;v),\\
 \Phi_a[F](v)
   &=P(v)G(a;v)+\sum_\alpha r_{\alpha1}G(\alpha,a;v)
                            -\Phi_{\varnothing}[P/(v-a)](v),\\
 \Phi_{a,b}[F](v)
   &=P(v)G(a,b;v)+\sum_\alpha r_{\alpha1}G(\alpha,a,b;v)
                            -\Phi_b[P/(v-a)](v).
 \label{eq:comp-supp-box-gpl-primitive}
\end{align}
Each recursive use shortens the word. The displayed formulas
therefore terminate after at most three algebraic steps.
Differentiating them gives
$\Phi_w'[F](v)=F(v)G(w;v)$, since the derivative of a nonempty
word removes its first letter and divides by its kernel denominator.
Thus these formulas describe closed GPL expressions of weight
at most three, not a residual parameter quadrature.

\paragraph{Endpoint expression for the box contribution.}
For each photon ordering $j\in\{s,u\}$, the preceding explicit
scalar formulas uniquely specify the finite set of rational
coefficients $F_{jw}(v;x)$ through
\begin{equation}
 \frac{8v}{(1-v^2)^2}
       \operatorname{Re}\mathcal T_{\mathrm{box},j}(x,y(v))
       =\sum_{|w|\leq2}F_{jw}(v;x)G(w;v).
\end{equation}
There are nine input words in the $s$ ordering and twenty-nine
in the $u$ ordering when products have been interleaved as above.
Apply \eqref{eq:comp-supp-box-gpl-primitive} to these rational
coefficients and define the resulting finite GPL sums by
\begin{equation}
 \mathcal F_j(v;x)=\sum_w\Phi_w[F_{jw}](v).
\end{equation}
All partial-fraction coefficients used here have the explicit
algebraic definition
\begin{equation}
 r_{aj}=\frac1{(m_a-j)!}
 \left.\left(\frac{\dd}{\dd v}\right)^{m_a-j}
                  [(v-a)^{m_a}F(v)]\right|_{v=a}.
\end{equation}
Thus a coefficient need not be expanded into a long quotient of
polynomials and square roots in order to define or evaluate the
closed answer.

The lower boundary is also determined by finite algebra.
Write the rational coefficient of a word in $\mathcal F_j$ as
$f_w(v)$, with pole order $p_w$ at zero. For a word that does not
end in zero, let $g_{w,n}$ be its Taylor coefficient of $v^n$.
The coefficients required at this boundary follow from
\begin{align}
 g_{\varnothing,n}&=\delta_{n0},\qquad g_{aw,0}=0,\\
 g_{aw,n}&=-\frac1n\sum_{k=0}^{n-1}
                   \frac{g_{w,k}}{a^{n-k}}\quad(a\ne0),\\
 g_{0w,n}&=\frac{g_{w,n}}n\quad(n\geq1).
\end{align}
Only coefficients up to $n=p_w$ are needed. The single logarithm
$G(0;v)=\ln v$ has zero constant term. Its divergent coefficient
cancels in the combined primitive, as required by the regular
original integrand at forward scattering. Define
\begin{equation}
 \mathcal L_j(x)=\sum_w\sum_{n=0}^{p_w}
 \frac{g_{w,n}}{(p_w-n)!}
 \left.\left(\frac{\dd}{\dd v}\right)^{p_w-n}
                    [v^{p_w}f_w(v)]\right|_{v=0}.
\end{equation}
This is an explicit finite sum of derivatives of rational functions.
It does not require an integration or a numerical endpoint limit.
The integrated box contribution to
$\sigma_1=\alpha_{\mathrm T}^3\mathcal C/m^2$ is
\begin{equation}
 \boxed{\mathcal C_{\mathrm{box}}(x;\lambda)
 =\frac1{2x^2}\sum_{j=s,u}
 \left[\mathcal F_j\!\left(\frac{x}{x+2};x\right)
                         -\mathcal L_j(x)\right].}
 \label{eq:comp-supp-box-angular-closed}
\end{equation}
The factor $1/(2x^2)$ follows from
$\dd c=2(x+1)\dd y/x^2$ and the interference normalization
$\mathcal C_{\mathrm{virt}}=(4s)^{-1}
\int\dd c\,\operatorname{Re}\mathcal T$.
For the preferred shorter construction let $\mathcal F_{u,\mathrm{rem}}$
be the same finite primitive operation applied only to the remainder
specified in Section~\ref{sec:comp-crossed-first}, and let
$\mathcal L_{u,\mathrm{rem}}$ be its lower boundary computed by the
same finite series rule. The full crossed boundary difference obeys
\begin{equation}
 \begin{aligned}
 \mathcal F_u(v_{\max};x)-\mathcal L_u(x)
 ={}&P(x,x)-P\!\left(x,\frac{x}{1+x}\right)\\
 &+\mathcal F_{u,\mathrm{rem}}(v_{\max};x)
       -\mathcal L_{u,\mathrm{rem}}(x).
 \end{aligned}
\end{equation}
Here $P(x,y)$ is the rational-subgroup primitive derived above.
The equality follows by differentiating the two constructions and using
the same original angular endpoints. It reduces the crossed angular
input from twenty-nine to seventeen words and avoids an internal
$\rho_u$ pole in this construction.

The expression still contains the infrared logarithm. It is combined
with the other virtual diagrams and the real contribution before
removing the photon mass.

\paragraph{Continuation through the apparent crossed-channel pole.}
For $x>1$, the positive roots obey $\eta<v_{\max}<Z$ because
\begin{equation}
 v_{\max}^2-\eta^2=\frac4{(x+2)^2(x+3)}>0,\qquad
 Z^2-v_{\max}^2=\frac{4x}{(x+4)(x+2)^2}>0.
\end{equation}
The original triangle requires the root $\eta-\ii0$.
An equivalent GPL evaluation starts at zero and deforms the path
into the upper half-plane before returning to $v_{\max}$.
In particular,
\begin{equation}
 G(\eta-\ii0;v_{\max})
 =\ln\!\left(\frac{v_{\max}}\eta-1\right)-\ii\pi.
\end{equation}
This follows by writing the weight-one integral as
$\log(v_{\max}-\eta+\ii0)-\log(-\eta+\ii0)$: the
first argument has phase zero and the second has phase $+\pi$.
Discarding this imaginary part term by term would change the
analytic continuation. Instead all words use the same deformed
path. Their imaginary parts cancel in the complete real
crossed-channel interference.
